\chapter{Conclusions and Suggestions}

\textcolor{red}{TODO: This chapter cannot be written until the experiment has been completed and the data analysed. The argument sketched out before the trial began was as follows, and should be rewritten against the actual results rather than kept as it stands: any differences in salt stress tolerance found between the species will be discussed in relation to the species-specific physiological responses reported by \textcite{Singh2022}; any mitigation of salt stress by biochar will be compared with the substrate effects described by \textcite{Joseph2021}; and the overall outcome will be related to the potential of biochar as a sustainable peat alternative under abiotic stress \textcite{Dumroese2018}.}

\section{Comparison of salt stress responses among mint species}

\section{Role of biochar in mitigating salt stress effects}

\section{Comparison with previous studies}

\section{Main conclusions}

\section{Practical implications and suggestions for indoor mint cultivation}

\section{Limitations of the study}

The principal limitation of this study is its limited replication. Each combination of growing medium and salinity was represented by a single pot within each species, and the leaves sampled from a pot are subsamples of that pot rather than independent replicates of the treatment. Differences observed between treatments may therefore partly reflect differences between individual pots, and the findings should be regarded as preliminary.

The unamended peat pots did not receive any salt, so the study contains no salted plants grown without biochar. The effect of biochar under saline conditions therefore cannot be fully separated from the effect of salinity itself. In addition, \textit{M. spicata} `Moroccan' had no biochar pots at 0~mM.

An aphid infestation occurred during the growing period. The insecticide was applied uniformly to every plant regardless of treatment, so plant protection itself is not confounded with the experimental treatments. Any differences in the severity of the infestation before treatment may nevertheless have affected growth independently of the growing medium and salinity applied.

The experiment was conducted over a single growing period at one location, and under greenhouse rather than field conditions. The results describe the behaviour of these species under the conditions tested and should not be extended to other environments, seasons or production systems without further work.

\section{Recommendations for future research}
